% ============================================
% 1. 问题重述
% ============================================

\section{问题重述}

\subsection{问题背景}

【宏观背景，约300字，需引用2-3篇参考文献】

随着\textbf{【领域背景】}的快速发展，\textbf{【具体问题】}日益受到关注。近年来，学者们围绕该问题开展了大量研究。\cite{ref1} 提出了\textbf{【已有方法1】}，该方法在\textbf{【场景】}中表现良好，但存在\textbf{【局限性1】}。\cite{ref2} 通过\textbf{【已有方法2】}对\textbf{【问题】}进行改进，然而\textbf{【局限性2】}仍未得到有效解决。\cite{ref3} 的研究表明，\textbf{【关键发现】}，这为解决该问题提供了新的思路。

当前，\textbf{【研究现状总结】}。因此，亟需一种\textbf{【本文要解决的问题】}的方法，以\textbf{【研究目标】}。

\subsection{问题提出}

本文所研究的问题来源于\textbf{【赛题来源】}。附件数据包含\textbf{【数据描述：文件类型、数据量、关键字段】}，涵盖了\textbf{【数据覆盖的时间范围/对象/维度】}等信息。

基于附件数据，本文需要解决以下问题：

\textbf{问题一：}【问题一的完整描述，包括输入、输出和约束条件】

\textbf{问题二：}【问题二的完整描述，包括输入、输出和约束条件】

\textbf{问题三：}【问题三的完整描述，包括输入、输出和约束条件】

\textbf{问题四：}【问题四的完整描述，包括输入、输出和约束条件】
